\documentclass[11pt,a4paper]{article}

% ---------------------------------------------------------
% PACKAGES
% ---------------------------------------------------------
\usepackage[utf8]{inputenc}
\usepackage[T1]{fontenc}
\usepackage{lmodern}
\usepackage{amsmath,amssymb,amsfonts,bm}
\usepackage{geometry}
\usepackage{graphicx}
\usepackage{hyperref}
\usepackage{authblk}
\usepackage{cite}
\usepackage{physics}
\usepackage{mathtools}
\usepackage{tensor}
\usepackage{enumitem}

\geometry{margin=2.4cm}

\title{\textbf{Companion XI (v49.1)}\\[0.2cm]
\textbf{Background Cosmology and Modified Friedmann Equations}\\[0.2cm]
\large Foundations of Relaxation-Driven Cyclic Cosmology (RDCC)}

\author[1]{Michael Lehmann}
\affil[1]{\small Independent Researcher, Relaxation-Driven Cyclic Cosmology (RDCC)\\
\texttt{mi.lehmann@gmx.de}}

\date{September 2026}

\begin{document}
\maketitle

\begin{abstract}
Relaxation-Driven Cyclic Cosmology (RDCC) describes the Universe as a globally
CPT-symmetric, bipartite quantum system whose observable deviations from
\(\Lambda\)CDM are governed by a single infrared quantity: the relaxation
parameter \(\alpha_{\mathrm{IR}}\). In v49.1 this parameter is promoted from a
phenomenological constant to a geometric quantity arising from a dimension-8
operator that modifies the large-scale structure of gravity. This Companion
develops the background cosmology of RDCC in its Phase-1-integrated v49.1
form. Starting from the global Einstein equation with the total stress-energy
tensor \(T^{(+)}_{\mu\nu}+T^{(-)}_{\mu\nu}\), we derive the effective
Friedmann equations as seen by a visible-sector observer, identify the
mirror sector as a non-clustering, diffuse component that behaves as
effective dark matter, and incorporate the geometric relaxation corrections
into the expansion history and distance--redshift relations. The resulting
background cosmology reproduces a \(\Lambda\)CDM-like expansion while
encoding the RDCC-specific scale dependence in the perturbation sector.
\end{abstract}

\tableofcontents
\newpage

\section{Introduction}

RDCC is built on the structural postulate that the global quantum state of
the Universe is a CPT eigenstate,
\begin{equation}
  \Theta\,|\Psi_{\mathrm{CPT}}\rangle = |\Psi_{\mathrm{CPT}}\rangle.
\end{equation}
The minimal realization of such a state requires a bipartite Hilbert-space
factorization,
\begin{equation}
  \mathcal{H} = \mathcal{H}_{(+)} \otimes \mathcal{H}_{(-)},
\end{equation}
where \(\mathcal{H}_{(+)}\) describes the visible sector and
\(\mathcal{H}_{(-)}\) its CPT-conjugate mirror sector.

Gravity is the unique interaction that couples to the sum of both sectors.
In v49.1 this global structure is retained, but the relaxation parameter is
reinterpreted as a geometric correction arising from a dimension-8 operator
that modifies the effective large-scale gravitational response. The purpose
of this Companion is to construct the background cosmology of RDCC in
v49.1: to derive the effective Friedmann equations, to identify the
mirror-sector contribution to the expansion history, and to incorporate the
geometric relaxation corrections into \(H(z)\) and cosmological distance
measures.

\section{Global Einstein equation and total stress-energy}

\subsection{Two-sector action and global coupling}

The global action of RDCC in v49.1 consists of the Einstein--Hilbert term,
the matter actions of the two sectors, and a higher-dimensional geometric
operator encoding relaxation:
\begin{equation}
  S
  =
  \frac{1}{16\pi G_N}
  \int d^{4}x\,\sqrt{-g}\, R
  +
  S_{(+)}[g_{\mu\nu},\Phi_{(+)}]
  +
  S_{(-)}[g_{\mu\nu},\Phi_{(-)}]
  +
  S_{\mathrm{rel}}[g_{\mu\nu}].
\end{equation}

The total stress-energy tensor sourcing the metric is
\begin{equation}
  T^{\mathrm{tot}}_{\mu\nu}
  =
  T^{(+)}_{\mu\nu}
  +
  T^{(-)}_{\mu\nu}
  +
  T^{\mathrm{rel}}_{\mu\nu},
\end{equation}
with
\begin{equation}
  T^{\mathrm{rel}}_{\mu\nu}
  =
  -\frac{2}{\sqrt{-g}}
  \frac{\delta S_{\mathrm{rel}}}{\delta g^{\mu\nu}}.
\end{equation}

Variation yields the global Einstein equation
\begin{equation}
  G_{\mu\nu}[g]
  =
  8\pi G_N
  \left(
    T^{(+)}_{\mu\nu}
    +
    T^{(-)}_{\mu\nu}
    +
    T^{\mathrm{rel}}_{\mu\nu}
  \right).
  \label{eq:globalEinstein}
\end{equation}

\subsection{Dimension-8 relaxation geometry}

In v49.1 the relaxation parameter arises from a dimension-8 operator of the
form
\begin{equation}
  \mathcal{L}_{\mathrm{rel}}
  =
  \frac{1}{M^{4}}\,
  C^{\mu\nu\alpha\beta}
  T^{(+)}_{\mu\nu} T^{(-)}_{\alpha\beta},
\end{equation}
where \(C^{\mu\nu\alpha\beta}\) is a \emph{dimensionless} geometric tensor.
The operator has mass dimension 4, as required for a Lagrangian density.
The effective modification of the Einstein tensor arises only after
variation with respect to the metric:
\begin{equation}
  G_{\mu\nu}
  \;\longrightarrow\;
  G_{\mu\nu} + \alpha_{\mathrm{IR}}\,\Delta G_{\mu\nu}.
\end{equation}

At the background level, \(\Delta G_{\mu\nu}\) reduces to a small correction
to the Friedmann equations.

\section{Sectoral projection and effective FRW background}

\subsection{FRW ansatz and sectoral energy densities}

We assume a spatially flat FRW metric,
\begin{equation}
  ds^{2} = -dt^{2} + a^{2}(t)\, d\mathbf{x}^{2},
\end{equation}
with Hubble parameter \(H=\dot{a}/a\). The stress-energy tensors of the two
sectors are modeled as perfect fluids,
\begin{align}
  T^{(+)}_{\mu\nu}
  &=
  (\rho_{(+)}+p_{(+)}) u_{\mu} u_{\nu} + p_{(+)} g_{\mu\nu},
  \\
  T^{(-)}_{\mu\nu}
  &=
  (\rho_{(-)}+p_{(-)}) u_{\mu} u_{\nu} + p_{(-)} g_{\mu\nu}.
\end{align}

Inserting these into Eq.~\eqref{eq:globalEinstein} yields the global
Friedmann equations,
\begin{align}
  3H^{2}
  &=
  8\pi G_N
  \left(
    \rho_{(+)} + \rho_{(-)} + \rho_{\mathrm{rel}}
  \right),
  \label{eq:globalFriedmann1}
  \\
  2\dot{H} + 3H^{2}
  &=
  -8\pi G_N
  \left(
    p_{(+)} + p_{(-)} + p_{\mathrm{rel}}
  \right).
  \label{eq:globalFriedmann2}
\end{align}

\subsection{Visible-sector projection}

A visible-sector observer measures only \(T^{(+)}_{\mu\nu}\) directly. The
effective Friedmann equations become
\begin{align}
  3H^{2}
  &=
  8\pi G_{\mathrm{eff}}\, \rho_{(+)} + \Lambda_{\mathrm{eff}},
  \label{eq:sectorFriedmann1}
  \\
  2\dot{H} + 3H^{2}
  &=
  -8\pi G_{\mathrm{eff}}\, p_{(+)}
  - 8\pi G_N \left( p_{(-)} + p_{\mathrm{rel}} \right),
  \label{eq:sectorFriedmann2}
\end{align}
with
\begin{equation}
  G_{\mathrm{eff}}
  =
  G_N
  \left(
    1
    +
    \frac{\rho_{(-)}}{\rho_{(+)}}
    +
    \frac{\rho_{\mathrm{rel}}}{\rho_{(+)}}
  \right).
  \label{eq:Geff}
\end{equation}
The effective coupling follows directly from the projection of the global
Einstein equation onto the visible sector. A detailed derivation is given
in Technical Note (v49.1): Effective Gravitational Coupling in RDCC. An
earlier inverted expression for \(G_{\mathrm{eff}}\) is incompatible with
Eq.~\eqref{eq:globalFriedmann1} and is superseded by Eq.~\eqref{eq:Geff}.
The \(\Omega\)-form of \(H(z)\) below is unchanged.

\section{Mirror sector as effective dark-matter component}

\subsection{Thermodynamic blockade}

The mirror sector contributes to the background expansion but does not
cluster on observable scales. Its energy density behaves as a diffuse,
non-clustering component:
\begin{equation}
  \rho_{\mathrm{DM,eff}} \equiv \rho_{(-)}.
\end{equation}

\section{Modified Friedmann equations and \texorpdfstring{\(H(z)\)}{H(z)}}

\subsection{Background components}

The global Friedmann equation becomes
\begin{equation}
  H^{2}(z)
  =
  H_0^{2}
  \left[
    \Omega_{\mathrm{b}}(1+z)^{3}
    +
    \Omega_{\mathrm{DM,eff}}(1+z)^{3}
    +
    \Omega_{\mathrm{rad}}(1+z)^{4}
    +
    \Omega_{\Lambda}^{\mathrm{eff}}
    +
    \delta_{\mathrm{rel}}(z)
  \right],
  \label{eq:Hofz}
\end{equation}
with
\begin{equation}
  \delta_{\mathrm{rel}}(z)
  =
  \alpha_{\mathrm{IR}}\, f_{\mathrm{rel}}(z).
\end{equation}

\section{Distance--redshift relations}

\subsection{Comoving distance}

\begin{equation}
  \chi(z)
  =
  \int_0^{z} \frac{dz'}{H(z')}.
\end{equation}

To leading order in \(\alpha_{\mathrm{IR}}\),
\begin{equation}
  H(z)
  =
  H_{\Lambda\mathrm{CDM}}(z)
  \bigl[
    1 + \alpha_{\mathrm{IR}}\, \epsilon(z)
  \bigr].
\end{equation}

\section{Conclusion}

This Companion establishes the background cosmology of RDCC in v49.1.
Gravity couples to the sum of both sectors, the mirror sector acts as a
diffuse geometric component, and the infrared relaxation parameter
introduces small, correlated corrections to the expansion history. The
sectoral coupling \(G_{\mathrm{eff}}\) is that of Companion I, as
consolidated in Technical Note (v49.1): Effective Gravitational Coupling
in RDCC.

Together with Companions I, III, IV, VI and VII, this Companion completes
the background foundation of RDCC and prepares the ground for the full
Boltzmann implementation and Euclid-oriented forecasting.

\end{document}